\section{Related Work}\label{sec:related_work}
Automated journalism is an interesting concept that involves the use of machine learning models and natural language processing. It can simply be defined as automatically generating news content from a variety of resources. It's been already claimed that the use of automation in news creation is changing journalism (\cite{guzman2019prioritizing}) and it will continue to do so. Here, we would like to discuss some of the prominent automated journalism applications that we find worth mentioning.
\subsection{Bloomberg’s Cyborg}
Bloomberg uses an application called Cyborg to help reporters with their financial reports about company earnings. This helps Bloomberg's journalists to create more complex financial articles (\cite{micklethwait2018future}). It creates news stories from a vast amount of financial articles and displays an overview.
\subsection{WordSmith by Associated Press (AP)}
AP uses Wordsmith to automatically generate earnings reports for various companies. Rather than manually extracting information from the financial reports quarterly released by the companies, it allows journalists or reporters to access well-structured and readable texts. In other words, WordSmith extracts the financial data and generates reports that sound fluent and plausible.
\subsection{Reuters News Tracer}
This tool helps Reuters journalists detect and verify newsworthy events on Twitter in real-time. In other words, it performs event detection using clustering to find more prominent events. Moreover, to determine which events are more newsworthy than others, it leverages a ranking algorithm. In addition, it's able to update the summaries based on the changing events in a time period, i.e., it has the ability to perform temporal summarization. Since working with social media data is relatively harder than structured reports or articles, we think Reuters News Tracer is an impressive real-time application.
\subsection{Quakebot by the Los Angeles Times}
The last application that draws our attention is Quakebot\footnote{\href{https://www.latimes.com/people/quakebot}{Quakebot, https://www.latimes.com/people/quakebot}}. It is a software application that reports the latest earthquakes using data from the U.S. Geological Survey. Its relation with automated journalism is it generates a draft article if the earthquake notices are above a certain threshold. It's reviewed by an editor and published if it's found newsworthy.